\section{Why Existing Approaches Are Not Interchangeable}
\label{sec:alternatives}
The relevant question is not whether another approach is secure in isolation,
but whether its authenticated principal and authorization decision reach the
executor for this interaction. Table~\ref{tab:alternatives} summarizes deployment
conditions rather than ranking unrelated products.

\begin{table*}[t]
\centering
\begin{tabularx}{\textwidth}{p{1.30in}XX}
\toprule
Approach & Where it can be sufficient & What remains at an ordinary voice boundary \\
\midrule
Authenticated business API & The agent and business expose compatible endpoints, with scoped authorization and suitable confirmation. & Often the preferred path. A separate call still needs correlation if its content or parties matter to execution. \\
OAuth, PAP, A2A, or MCP & Discovery, agent identity, account access, and tool interactions through supported authenticated transports. & A valid web session is not automatically a proof about an unrelated caller. A transport binding can connect them. \\
WebRTC / SRTP & Media confidentiality and integrity between the actual cryptographic media endpoints. & Media keys alone do not specify customer delegation; a gateway terminating media may be a different trust boundary. \\
STIR / SIP identity & Asserted telephone identity under the signaling trust model. & A telephone number is not a scoped, single-use customer authorization for an application operation. \\
Synthetic-voice detection or biometrics & Risk signals, triage, and identification within a deployed detection model. & Recognizing a voice does not demonstrate possession of an agent key or identify the exact customer grant. \\
Spoken codes or user callbacks & Explicit user participation with a correctly designed transaction-binding flow. & Codes can be relayed or disclosed; exact intent and concurrent-call binding need a protocol and user interface. \\
TLS over audio / QUIC over audio & Protected machine requests through a mutually supported audio path. & Requires both endpoints, a decodable modem, trusted authorization, and a separately justified enrollment policy. \\
\bottomrule
\end{tabularx}
\caption{Complementary approaches and the assumptions under which they address the voice-to-executor boundary. An available authenticated API can remove the need for an audio control channel.}
\label{tab:alternatives}
\end{table*}

\paragraph{Web authorization is a foundation.}
OAuth certificate-bound tokens and DPoP already bind credentials to possession
of a key~\cite{rfc8705,rfc9449}. PAP describes personal-agent discovery, sessions,
account access, and conversations; its current public draft is 0.3~\cite{pap03}.
MCP defines HTTP authorization, and A2A provides agent interaction and security
mechanisms~\cite{mcp,a2a}. We do not claim these protocols lack security. The
additional question is how a particular audio endpoint demonstrates possession
of the relevant key and presents permission to its executor. Our prototype
implements a custom bridge for the repository's earlier PAP draft 0.1 profile;
it is not evidence of current-draft conformance or external interoperability.

An HTTPS grant sent to a receiving application can be entirely adequate when
the receiving application securely correlates that grant with the call. Such
an approach is an important baseline, not a design to exclude. A spoken call ID
alone is a weak correlation primitive if an attacker can relay it; adding
authenticated channel evidence changes the comparison. We plan to compare
these concrete designs, including their operational costs.

\paragraph{Media encryption and signaling identity answer different questions.}
WebRTC security protects media between its selected endpoints~\cite{rfc8826}.
It does not itself express the customer's permitted account operation. SIP
authenticated identity addresses assertions about originating identity in
signaling~\cite{rfc8224}; it does not replace an application grant. Inner
cryptography is useful if media must traverse an untrusted terminator while
the business executor and agent retain their keys. If media and authorization
already terminate at the intended trusted applications, an additional audio
tunnel may be redundant.

\paragraph{Detection remains useful.}
Detection can operate on callers that do not participate in a new protocol and
can identify suspicious behavior even after authentication. Conversely, a
known voice can be shared by many users, and a valid certificate can belong to
a malicious agent. Treating detection as a risk input and explicit grants as
execution constraints preserves the value of both. \sys{} does not establish
a measured improvement over a commercial detector.

\section{What Makes the Problem Non-Trivial}
\label{sec:challenges}
\paragraph{The modality is not a packet socket.}
An audio path introduces serialization delay, codec transformations, framing
loss, timing drift, echo control, and speech interference. A link that preserves
human intelligibility need not preserve modem symbols. A certificate chain
small on an IP network can occupy seconds of a low-rate call, interacting with
loss timers and grant expiry.

\paragraph{Loss recovery must be placed deliberately.}
TLS records need reliable ordered bytes; the existing profile supplies them.
QUIC supplies recovery for its own handshake and reliable streams, allowing
the audio link to discard damaged datagrams. However, QUIC requires support for
1200-byte Initial datagrams~\cite{rfc9000}. A 256-byte audio payload therefore
requires fragmentation, without converting incomplete datagrams into a global
ordered backlog. Under an illustrative independent fragment-erasure probability
$p$, a five-fragment packet survives with probability $(1-p)^5$. At $p=0.01$
this predicts about $4.9\%$ packet erasure. This is an analytical illustration,
not a channel measurement; burst loss will violate the independence assumption.

\paragraph{Authority needs more than a secure handshake.}
A certificate identifies a key under an issuer's policy. A customer grant
restricts what that key can do. A live challenge establishes freshness. An
argument digest can bind one exact proposal. None implies that the surrounding
speech is truthful or that a model's interpretation matches the user's intent.
The executor must keep these assertions separate and validate them at the
right point in its transaction.

\paragraph{Conversation and reliable execution have different semantics.}
Two modem waveforms must not interleave, and speech should not overwrite a
protected exchange. Even a perfectly recovered request can execute once while
its response is lost. A retry policy cannot infer whether the application
committed a write. Finally, finishing an application receive does not mean its
delayed transport acknowledgment has finished playing. Teardown must preserve
that final audio, especially on the slower modem.
